\documentclass[11pt,a4paper]{article}
\usepackage[margin=2.7cm]{geometry}
\usepackage[T1]{fontenc}
\usepackage[utf8]{inputenc}
\usepackage[ngerman]{babel}
\usepackage{lmodern}
\usepackage{amsmath,amssymb,amsthm,mathtools}
\usepackage{booktabs,array}
\usepackage{microtype}
\usepackage[hidelinks]{hyperref}
\usepackage{xcolor}
\usepackage{enumitem}
\usepackage{bm}

\newtheorem{theorem}{Satz}[section]
\newtheorem{proposition}[theorem]{Proposition}
\newtheorem{lemma}[theorem]{Lemma}
\newtheorem{corollary}[theorem]{Korollar}
\theoremstyle{definition}
\newtheorem{definition}[theorem]{Definition}
\theoremstyle{remark}
\newtheorem{remark}[theorem]{Bemerkung}

\newcommand{\PP}{\mathbb P}
\newcommand{\NN}{\mathbb N}
\newcommand{\OO}{\mathbb O}
\newcommand{\HH}{\mathbb H}
\newcommand{\RR}{\mathbb R}
\newcommand{\Sym}{\operatorname{Sym}}
\newcommand{\Aut}{\operatorname{Aut}}
\newcommand{\Span}{\operatorname{span}}

\title{\textbf{Arithmetische Besetzungsspektren und die erste oktonionische Nichtassoziativitätsschwelle}}
\author{EABC Research Note}
\date{September 2026}

\begin{document}
\maketitle

\begin{abstract}
Wir führen eine kanonische Besetzungsraum-Darstellung der natürlichen Zahlen ein, die ausschließlich auf der eindeutigen Primfaktorzerlegung beruht. Für
\[
N=\prod_{p\in\PP}p^{\nu_p}
\]
ist das zugehörige Faktorwort
\[
W(N)=\sum_p\nu_p e_p,
\]
mit den arithmetischen Besetzungszahlen
\[
\Omega(N)=\sum_p\nu_p,
\qquad
\omega(N)=\#\{p:\nu_p>0\}.
\]
Der logarithmische diagonale Operator besitzt den exakten Spektralwert
\[
E_\times(N)=\sum_p\nu_p\log p=\log N.
\]
Primzahlpotenzen sind damit Mehrfachbesetzungen eines atomaren Primmodus, während Produkte verschiedener Primzahlen echte Mehrmodenzustände bilden. Wir zeigen, dass sich im Sektor $\Omega=3$ die drei Besetzungstypen $p^3$, $p^2q$ und $pqr$ zu einer exakten Hierarchie ordnen: Aufgrund der Alternativität und des Satzes von Artin ist genuine Drei-Erzeuger-Nichtassoziativität für die ersten beiden Typen unmöglich und wird erstmals im quadratfreien Typ $pqr$ möglich. Anschließend trennen wir das euklidische Tetraedervolumen von der $G_2$-Kalibrierung und vom oktonionischen Assoziatordefekt,
\[
V_{\rm tet}=\frac16\sqrt{\det G},\qquad
\varphi(u,v,w)=\langle uv,w\rangle,\qquad
\mathfrak t=\frac14\|[u,v,w]\|^2=\det G-\varphi^2.
\]
Eine dynamische Prime-Orbit-Realisierung oder physikalische Interpretation von $E_\times$ wird nicht vorausgesetzt.
\end{abstract}

\section{Faktorwort-Normalform}
Sei
\[
\mathfrak M=\bigoplus_{p\in\PP}\NN_0 e_p
\]
das freie kommutative Monoid über den Primzahlbezeichnungen $e_p$. Aufgrund der eindeutigen Primfaktorzerlegung besitzt jedes $N\ge1$ eine eindeutige Darstellung
\[
N=\prod_{p\in\PP}p^{\nu_p}
\]
mit nur endlich vielen von Null verschiedenen $\nu_p$ und damit ein kanonisches Faktorwort
\begin{equation}
W(N)=\sum_p\nu_p e_p.
\end{equation}
Zusammengesetzte Teilprodukte erzeugen keine neuen Generatoren. Somit gilt
\[
W(pq)=e_p+e_q,\qquad
W(p^2q)=2e_p+e_q,\qquad
W(pqr)=e_p+e_q+e_r.
\]

\begin{definition}[Arithmetische Besetzungszahlen]
Wir definieren
\[
\Omega(N)=\sum_p\nu_p,
\qquad
\omega(N)=\#\{p:\nu_p>0\}.
\]
Dabei zählt $\Omega$ die Primfaktoren mit Multiplizität, während $\omega$ die Zahl der verschiedenen besetzten Primmoden zählt.
\end{definition}

\section{Logarithmisches Besetzungsspektrum}
Wir betrachten die Hilbertbasis $\{|N\rangle:N\in\NN\}$, indiziert durch Faktorwörter. Definiere den diagonalen Operator
\begin{equation}
H_\times|N\rangle
=\left(\sum_p\nu_p\log p\right)|N\rangle.
\end{equation}

\begin{proposition}[Exakte arithmetische Energie]
Für jedes $N\in\NN$ gilt
\[
E_\times(N)=\log N.
\]
Ferner gilt für alle $M,N\in\NN$
\[
E_\times(MN)=E_\times(M)+E_\times(N).
\]
\end{proposition}

\begin{proof}
Dies folgt unmittelbar aus
\[
\log\left(\prod_pp^{\nu_p}\right)=\sum_p\nu_p\log p
\]
und $\log(MN)=\log M+\log N$.
\end{proof}

Die atomaren Primmodus-Energien sind somit
\[
\varepsilon_p=\log p,
\]
und ein allgemeiner Besetzungszustand besitzt den freien additiven Spektralwert
\[
E_\times(N)=\sum_p\nu_p\varepsilon_p.
\]
Dies ist eine arithmetische Spektralaussage und noch keine physikalische Energiebehauptung.

\section{Primale, semiprimale und dreifache Besetzungssektoren}
Für $\Omega=1$ existieren ausschließlich Primzustände
\[
p:\quad W(p)=e_p,\qquad E_\times(p)=\log p.
\]
Für $\Omega=2$ existieren genau zwei Besetzungstypen:
\[
p^2:\quad W=2e_p,\quad \omega=1,\quad E_\times=2\log p,
\]
\[
pq\ (p\ne q):\quad W=e_p+e_q,\quad \omega=2,\quad E_\times=\log p+\log q.
\]
Damit ist $p^2$ eine Mehrfachbesetzung eines einzigen Modus, während $pq$ ein echter Zweimodenzustand ist.

Für $\Omega=3$ liefert die eindeutige Primfaktorzerlegung exakt die Partitionstypen $3$, $2+1$ und $1+1+1$:
\begin{center}
\begin{tabular}{@{}llll@{}}
\toprule
Arithmetischer Typ & Faktorwort & $\omega(N)$ & $E_\times(N)$ \\
\midrule
$p^3$ & $3e_p$ & $1$ & $3\log p$ \\
$p^2q$, $p\ne q$ & $2e_p+e_q$ & $2$ & $2\log p+\log q$ \\
$pqr$, paarweise verschieden & $e_p+e_q+e_r$ & $3$ & $\log p+\log q+\log r$ \\
\bottomrule
\end{tabular}
\end{center}

In bosonischer Besetzungsnotation entsprechen diese formal
\[
|3_p\rangle,\qquad |2_p,1_q\rangle,\qquad |1_p,1_q,1_r\rangle
\]
in $\Sym^3\mathcal H_1$. Diese symmetrische Besetzungsstruktur darf nicht mit Exteriorpotenzen aus Fredholm-Determinantenentwicklungen identifiziert werden.

\section{Oktonionische Realisierung besetzter Moden}
Wir nehmen an, dass jedem verschiedenen besetzten Primmodus eine feste imaginäre oktonionische Richtung zugeordnet wird,
\[
p\longmapsto u_p\in\operatorname{Im}\OO.
\]
Eine Kanonizität dieser Zuordnung wird nicht vorausgesetzt.

Für $x,y,z\in\OO$ definieren wir den Assoziator
\[
[x,y,z]=(xy)z-x(yz).
\]
Die Oktonionen sind alternativ, also
\[
[x,x,y]=[y,x,x]=0.
\]
Darüber hinaus besagt der Satz von Artin, dass jede von zwei Elementen einer alternativen Algebra erzeugte Unteralgebra assoziativ ist.

\begin{theorem}[Besetzungs--Assoziativitäts-Schwelle]\label{thm:threshold}
Sei $N$ mit $\Omega(N)=3$ gegeben, und jedem verschiedenen besetzten Primmodus sei ein oktonionischer Generator zugeordnet. Dann gilt:
\begin{enumerate}[label=(\roman*)]
\item Für $N=p^3$ verschwindet jeder lokale Dreiwort-Assoziator.
\item Für $N=p^2q$ mit $p\ne q$ verschwindet jeder Assoziator, der aus den beiden besetzten Generatoren $u_p,u_q$ gebildet wird.
\item Für $N=pqr$ mit paarweise verschiedenen $p,q,r$ wird ein nichtverschwindender Assoziator $[u_p,u_q,u_r]$ erstmals möglich.
\end{enumerate}
Folglich gilt
\[
\omega(N)\le2
\quad\Longrightarrow\quad
\text{genuine Drei-Erzeuger-Nichtassoziativität ist unmöglich},
\]
während $\omega(N)=3$ die erste Besetzungsschwelle ist, an der sie auftreten kann.
\end{theorem}

\begin{proof}
Für $p^3$ sind alle drei Einträge derselbe Generator, daher folgt aus der Alternativität $[u_p,u_p,u_p]=0$. Für $p^2q$ treten höchstens die beiden Generatoren $u_p,u_q$ auf. Nach dem Satz von Artin ist die von ihnen erzeugte Unteralgebra assoziativ, weshalb jeder aus diesen Generatoren gebildete Dreiwort-Assoziator verschwindet. Für $pqr$ treten drei verschiedene Generatoren auf; der Zwei-Erzeuger-Satz von Artin erzwingt dann keine Assoziativität mehr. Nichtverschwindende Assoziatoren sind daher möglich, aber nicht notwendig.
\end{proof}

\begin{remark}
Die Quadratfreiheit impliziert keine Nichtassoziativität. Drei verschiedene oktonionische Generatoren können weiterhin in einer gemeinsamen assoziativen quaternionischen Unteralgebra liegen. Der Satz identifiziert den ersten \emph{möglichen} nichtassoziativen Besetzungstyp, nicht ein hinreichendes arithmetisches Kriterium.
\end{remark}

\section{$G_2$-Geometrie der Dreimodenkonfiguration}
Seien $u,v,w\in\operatorname{Im}\OO\cong\RR^7$. Wir bezeichnen mit
\[
G=(\langle u_i,u_j\rangle)_{i,j=1}^3
\]
ihre Gram-Matrix. Das euklidische Volumen des Simplex $\{0,u,v,w\}$ lautet
\begin{equation}
V_{\rm tet}(u,v,w)=\frac16\sqrt{\det G}.
\end{equation}
Dies ist eine rein metrische Invariante und darf nicht mit dem Assoziatordefekt verwechselt werden.

Die standardmäßige assoziative $G_2$-Dreiform ist bis auf die Orientierungskonvention
\begin{equation}
\varphi(u,v,w)=\langle uv,w\rangle.
\end{equation}
Definiere den Nichtassoziativitätsdefekt
\begin{equation}
\mathfrak t(u,v,w)=\frac14\|[u,v,w]\|^2.
\end{equation}
Für imaginäre Oktonionen gilt die Identität
\begin{equation}\label{eq:g2identity}
\boxed{\mathfrak t=\det G-\varphi^2.}
\end{equation}
Äquivalent dazu gilt
\begin{equation}
\boxed{\mathfrak t=36V_{\rm tet}^2-\varphi^2.}
\end{equation}
Somit sind $V_{\rm tet}$, $\varphi$ und $\mathfrak t$ miteinander verknüpfte Observablen und keine drei unabhängigen Invarianten.

Für ein orthonormales Tripel gilt $G=I_3$ und damit
\[
V_{\rm tet}=\frac16,
\qquad
\mathfrak t=1-\varphi^2,
\qquad
0\le\mathfrak t\le1.
\]
Allgemeine orthonormale $G_2$-Tripel müssen nicht binär sein: Zwischenwerte von $|\varphi|$ und $\mathfrak t$ sind möglich.

\section{Diskreter Fano-Basissektor}
Seien $e_1,\dots,e_7$ die sieben imaginären Basiseinheiten einer festen Oktonionen-Multiplikationstafel. Es existieren
\[
\binom73=35
\]
unordnete Tripel verschiedener Basisrichtungen. Genau sieben davon sind Fano-Linien-Tripel und achtundzwanzig sind Nicht-Fano-Tripel.

Jedes Basistripel ist orthonormal. Daher erfüllen alle fünfunddreißig zugehörigen euklidischen Tetraeder
\[
V_{\rm tet}=\frac16.
\]
Die $G_2$-Multiplikation trennt sie dennoch:
\begin{center}
\begin{tabular}{@{}lcccc@{}}
\toprule
Typ & Anzahl & $\det G$ & $|\varphi|$ & $\mathfrak t$ \\
\midrule
Fano-Linie & $7$ & $1$ & $1$ & $0$ \\
Nicht-Fano-Basistripel & $28$ & $1$ & $0$ & $1$ \\
\bottomrule
\end{tabular}
\end{center}
Damit sind die fünfunddreißig Basistetraeder euklidisch volumengleich, aber oktonionisch nicht äquivalent. Der Unterschied liegt in der kalibrierten Multiplikationsstruktur und nicht im gewöhnlichen metrischen Volumen.

\section{Konfigurationsinvarianten versus arithmetische Invarianten}
Für ein festes Tripel $(u_p,u_q,u_r)$ sind die Größen
\[
G,\qquad V_{\rm tet},\qquad\varphi,\qquad\mathfrak t
\]
mathematisch wohldefiniert. Sie sind daher Konfigurationsinvarianten.

Sie sind jedoch noch keine Invarianten des bloßen Primzahltripels $(p,q,r)$. Um eine solche arithmetische Invariante zu erhalten, benötigt man eine kanonische oder quotientstabile Zuordnung
\[
p\longmapsto u_p
\]
oder ein invariantes Mittelungsverfahren über eine mathematisch spezifizierte Darstellungsfaser. Die Wahl einer einzelnen Vierquadratedarstellung
\[
p=a^2+b^2+c^2+d^2
\]
und eine willkürliche Einbettung dieser Koordinaten in $\operatorname{Im}\OO$ liefern eine solche Kanonizität nicht.

Insbesondere löst auch die Quotientierung eines einzelnen Einheitsvektors nach der vollen Gruppe $G_2=\Aut(\OO)$ das Problem nicht: $G_2$ wirkt transitiv auf der Einheitssphäre $S^6\subset\operatorname{Im}\OO$. Relevante $G_2$-Invarianten entstehen erst für Mehrvektorkonfigurationen.

\section{Primzahlpotenzen, Mehrmodenzustände und mögliche Orbit-Realisierungen}
Die exakte arithmetische Identität
\[
E_\times(p^k)=k\log p
\]
zeigt Primzahlpotenzen als Besetzungsleiter
\[
|1_p\rangle,\ |2_p\rangle,\ |3_p\rangle,\dots
\]
entlang eines einzigen Primmodus. Gemischte Produkte aktivieren mehrere Modi.

Die Formel ähnelt dem Wiederholungsgesetz
\[
\ell(\gamma^k)=k\ell(\gamma)
\]
für periodische Orbits. Die Existenz einer Dynamik mit
\[
p\longmapsto\gamma_p,\qquad \ell(\gamma_p)=\log p
\]
ist jedoch ein separates Realisierungsproblem und wird hier nicht vorausgesetzt. Insbesondere darf die arithmetische Gleichheit $E_\times(p)=\log p$ ohne zusätzlichen Beweis nicht mit einer physikalisch oder dynamisch realisierten primitiven Periode identifiziert werden.

Ebenso dürfen bosonische Besetzungssektoren
\[
\Sym^k\mathcal H_1
\]
nicht mit Fredholm-Exteriorpotenzen
\[
\Lambda^k\mathcal H
\]
verwechselt werden. Jede Brücke zwischen diesen Konstruktionen benötigt einen unabhängigen Satz.

\section{Epistemischer Status und offene Gates}
Die folgenden Aussagen sind innerhalb des algebraischen Rahmens etabliert:
\begin{itemize}
\item das kanonische Faktorwort $W(N)=\sum_p\nu_pe_p$;
\item die arithmetischen Besetzungsdaten $\Omega(N)$ und $\omega(N)$;
\item das exakte logarithmische Spektrum $E_\times(N)=\log N$;
\item die vollständige $\Omega=3$-Zerlegung in $p^3$, $p^2q$ und $pqr$;
\item Satz~\ref{thm:threshold};
\item für ein festes oktonionisches Tripel die geometrischen Größen $G$, $V_{\rm tet}$, $\varphi$ und $\mathfrak t$ sowie die Identität~\eqref{eq:g2identity};
\item die diskrete Fano-Basisklassifikation $35=7+28$.
\end{itemize}

Die folgenden Punkte bleiben offene Realisierungsprobleme:
\begin{enumerate}
\item eine kanonische arithmetische Zuordnung $p\mapsto u_p$;
\item Darstellungsunabhängigkeit der oktonionischen Konfigurationsinvarianten;
\item eine periodische Orbit-Realisierung mit primitiven Perioden $\log p$;
\item eine physikalische Interpretation von $E_\times$ als messbare Energie;
\item jede Identifikation zwischen Faktorwort-Fockbesetzung und Fredholm-Exteriorpotenz-Kombinatorik.
\end{enumerate}

\section{Schlussfolgerung}
Die eindeutige Primfaktorzerlegung bestimmt einen kanonischen Besetzungsvektor, und die logarithmische Auswertung überführt die multiplikative Arithmetik in ein exakt additives Spektrum:
\[
\boxed{
N=\prod_pp^{\nu_p}
\quad\longmapsto\quad
W(N)=\sum_p\nu_pe_p,
\qquad
E_\times(N)=\sum_p\nu_p\log p=\log N.
}
\]
Primzahlpotenzen kodieren Mehrfachbesetzungen eines atomaren Modus, während gemischte Produkte echte Modenvielfalt kodieren. Bei Besetzungszahl drei erzwingen Alternativität und der Satz von Artin, dass die Sektoren $p^3$ und $p^2q$ assoziativ bleiben, und identifizieren den quadratfreien Sektor $pqr$ als den ersten, in dem ein genuiner oktonionischer Assoziator auftreten kann.

Daraus ergibt sich das Strukturprinzip
\[
\boxed{
\text{arithmetische Modenkomplexität}
\;\Longrightarrow\;
\text{erste mögliche algebraische Nichtassoziativitätsschwelle}.
}
\]
Diese Schlussfolgerung ist rein algebraisch. Die verbleibende Aufgabe ist die arithmetische Verankerung der oktonionischen Konfiguration, nicht die interne $G_2$-Geometrie selbst.

\end{document}
